% Unit 046 English translation of chapter3.tex lines 2936--3200.
% Source: Methods of Algebra, Volume 2, commit
% 9a5803ff77dd3257484cb177f851a73770a59dd3 (CC BY 4.0).
% stable-unit-id: o014.aljabr2.chapter3.k-injectives
% Coverage: Section 3.15 in full.

% segment-id: o014.li2.chapter3.k-injectives.h001
\section{K-injective and K-projective complexes}\label{sec:K-injectives}
\hypertarget{o014.aljabr2.chapter3.k-injectives}{ }

% segment-id: o014.li2.chapter3.k-injectives.p001
Injective and projective resolutions of complexes supplied the initial
definition of derived functors in \S\ref{sec:derived-primer}. That construction
requires bounded-below or bounded-above complexes, whereas applications
involving unbounded complexes and their derived categories require a broader
class of resolutions. The theory of K-injective and K-projective resolutions
originates with \cite{Spa88}. This section provides some preparation for that
theory; the arguments follow \cite{stacks}, and the reader may also consult
\cite[Chapter 14]{KS06} or \cite{BN93}.

% segment-id: o014.li2.chapter3.k-injectives.de001
\begin{definition}[N.\ Spaltenstein]\label{def:K-resolutions}
	\index{complex!K-injective, K-projective}
	\index{resolution!K-injective, K-projective}
	\index{enough K-injective complexes}
	\index{enough K-projective complexes}
	Let $\mathcal{A}$ be an abelian category.
% segment-id: o014.li2.chapter3.k-injectives.l001
	\begin{itemize}
% segment-id: o014.li2.chapter3.k-injectives.i001
		\item A complex $I\in\Obj(\cate{C}(\mathcal{A}))$ is called a
		\emph{K-injective complex}\footnote{A perhaps more natural term
		would be homotopically injective complex, and similarly
		homotopically projective complex.} if, for every acyclic complex
		$X\in\Obj(\cate{C}(\mathcal{A}))$, the complex
		$\Hom^\bullet(X,I)$ is also acyclic.

		If $f:A\to I$ is a quasi-isomorphism in
		$\cate{C}(\mathcal{A})$ and $I$ is K-injective, then $f$ is called
		a \emph{K-injective resolution} of $A$.
% segment-id: o014.li2.chapter3.k-injectives.i002
		\item A complex $P\in\Obj(\cate{C}(\mathcal{A}))$ is called a
		\emph{K-projective complex} if, for every acyclic complex
		$X\in\Obj(\cate{C}(\mathcal{A}))$, the complex
		$\Hom^\bullet(P,X)$ is also acyclic.

		If $f:P\to A$ is a quasi-isomorphism in
		$\cate{C}(\mathcal{A})$ and $P$ is K-projective, then $f$ is called
		a \emph{K-projective resolution} of $A$.
	\end{itemize}

	If every object of $\cate{C}(\mathcal{A})$ has a K-injective resolution
	(respectively, a K-projective resolution), then $\mathcal{A}$ is said to
	have \emph{enough K-injective complexes} (respectively,
	\emph{enough K-projective complexes}).
\end{definition}

% segment-id: o014.li2.chapter3.k-injectives.p002
The two sets of definitions are plainly dual. In practice, the following
criterion is often easier to use.

% segment-id: o014.li2.chapter3.k-injectives.le001
\begin{lemma}\label{prop:K-injective-criterion}
	A complex $I$ (respectively, $P$) is K-injective (respectively,
	K-projective) if and only if, for every acyclic complex $X$,
	$\Hom_{\cate{K}(\mathcal{A})}(X,I)=0$ (respectively,
	$\Hom_{\cate{K}(\mathcal{A})}(P,X)=0$).
\end{lemma}

% segment-id: o014.li2.chapter3.k-injectives.pf001
\begin{proof}
	Let $n\in\Z$. Lemma~\ref{prop:Hom-cplx-d-translate} gives
	\[
	\Hm^n \Hom^\bullet(X, I) \simeq
	\Hom_{\cate{K}(\mathcal{A})}(X[-n], I), \quad
	\Hm^n \Hom^\bullet(P, X) \simeq
	\Hom_{\cate{K}(\mathcal{A})}(P, X[n]);
	\]
	but $X$ is acyclic if and only if $X[-n]$ is acyclic, and this is also
	equivalent to $X[n]$ being acyclic.
\end{proof}

% segment-id: o014.li2.chapter3.k-injectives.ex001
\begin{example}\label{eg:bdd-K-resolution}
	Lemma~\ref{prop:resolution-homotopy-aux} can immediately be restated as
	follows.
% segment-id: o014.li2.chapter3.k-injectives.l002
	\begin{compactitem}
% segment-id: o014.li2.chapter3.k-injectives.i003
		\item If $I\in\Obj(\cate{C}^+(\mathcal{A}))$ consists of injective
		objects, then $I$ is K-injective.
% segment-id: o014.li2.chapter3.k-injectives.i004
		\item If $P\in\Obj(\cate{C}^-(\mathcal{A}))$ consists of projective
		objects, then $P$ is K-projective.
	\end{compactitem}
	Thus an injective (respectively, projective) resolution is a special case
	of a K-injective (respectively, K-projective) resolution.
\end{example}

% segment-id: o014.li2.chapter3.k-injectives.ex002
\begin{example}[A.\ Dold]
	The one-sided boundedness condition in Example~\ref{eg:bdd-K-resolution}
	is essential. For example, take
	$\mathcal{A}=\Z/4\Z\dcate{Mod}$ and consider the acyclic complex $Q$
	in it:
	\[
	\cdots \xrightarrow{2} \Z/4\Z \xrightarrow{2} \Z/4\Z
	\xrightarrow{2} \Z/4\Z \xrightarrow{2} \cdots.
	\]
	Every term is a free module. On the other hand, the termwise tensor
	product $Q\dotimes{\Z/4\Z}\Z/2\Z$ is the complex
	\[
	\cdots \xrightarrow{0} \Z/2\Z \xrightarrow{0} \Z/2\Z
	\xrightarrow{0} \Z/2\Z \xrightarrow{0} \cdots.
	\]
	Therefore $Q$ is not K-projective. Indeed, if $Q$ were K-projective,
	then $\identity_Q$ would be homotopic to $0$, and hence so would
	$\identity_{Q\otimes\Z/2\Z}$, contradicting
	$\Hm^n(Q\dotimes{\Z/4\Z}\Z/2\Z)=\Z/2\Z$.

	From another point of view, this example also shows that an unbounded
	complex consisting of projective objects need not be an appropriate
	resolution: $0\to Q$ and $0\to0$ are both quasi-isomorphisms, but the
	preceding paragraph shows that $\identity_Q$ is not homotopic to $0$.
	Thus $Q$ and $0$ are not isomorphic in $\cate{K}(\mathcal{A})$, and
	Theorem~\ref{prop:resolution-homotopy} cannot simply be extended to the
	unbounded case.
\end{example}

% segment-id: o014.li2.chapter3.k-injectives.p003
The following result is the natural generalization of
Theorem~\ref{prop:resolution-homotopy}; its proof is likewise postponed until
the discussion following
Theorem~\sourcecrossref{prop:cplx-triangulated}{in \S~4.4}.

% segment-id: o014.li2.chapter3.k-injectives.th001
\begin{theorem}\label{prop:K-resolution-homotopy}
	Let $\alpha:X\to Y$ be a quasi-isomorphism in
	$\cate{C}(\mathcal{A})$.
% segment-id: o014.li2.chapter3.k-injectives.l003
	\begin{itemize}
% segment-id: o014.li2.chapter3.k-injectives.i005
		\item Given a morphism $\gamma:X\to I$, where $I$ is a
		K-injective complex, there is a commutative diagram in
		$\cate{K}(\mathcal{A})$
% segment-id: o014.li2.chapter3.k-injectives.g001
		\[\begin{tikzcd}[row sep=tiny]
			& Y \arrow[dd, "\beta"] \\
			X \arrow[ru, "\alpha"] \arrow[rd, "\gamma"'] & \\
			& I
		\end{tikzcd}\]
% segment-id: o014.li2.chapter3.k-injectives.i006
		\item Given a morphism $\gamma:P\to Y$, where $P$ is a
		K-projective complex, there is a commutative diagram in
		$\cate{K}(\mathcal{A})$
% segment-id: o014.li2.chapter3.k-injectives.g002
		\[\begin{tikzcd}[row sep=tiny]
			& X \arrow[ld, "\alpha"'] \\
			Y & \\
			& P \arrow[lu, "\gamma"] \arrow[uu, "\beta"']
		\end{tikzcd}\]
	\end{itemize}
	In both cases, $\beta$ is unique as a morphism in
	$\cate{K}(\mathcal{A})$.
\end{theorem}

% segment-id: o014.li2.chapter3.k-injectives.pr001
\begin{proposition}\label{prop:K-sigma-duality}
	The equivalence $\sigma$ of
	Definition--Proposition~\ref{def:sigma} has the following property:
	$X\in\Obj(\cate{C}(\mathcal{A}^{\opp}))$ is K-injective (respectively,
	K-projective) if and only if
	$\sigma X\in\Obj(\cate{C}(\mathcal{A})^{\opp})
	=\Obj(\cate{C}(\mathcal{A}))$ is K-projective (respectively,
	K-injective).
\end{proposition}

% segment-id: o014.li2.chapter3.k-injectives.pf002
\begin{proof}
	The functor $\sigma$ preserves acyclic complexes and induces
	$\cate{K}(\mathcal{A}^{\opp})\rightiso
	\cate{K}(\mathcal{A})^{\opp}$
	(Proposition~\ref{prop:sigma-homotopy}); hence $\sigma$ preserves the
	criterion of Lemma~\ref{prop:K-injective-criterion}.
\end{proof}

% segment-id: o014.li2.chapter3.k-injectives.p004
Next consider adjoint functors. Proposition~\ref{prop:adjoint-injective-projective}
has a corresponding version in this setting.

% segment-id: o014.li2.chapter3.k-injectives.pr002
\begin{proposition}\label{prop:adjoint-injective-projectives-K}
	Consider a pair of functors between abelian categories
% segment-id: o014.li2.chapter3.k-injectives.g003
	$\begin{tikzcd}
		\mathcal{A} \arrow[r, shift left, "F"] &
		\mathcal{B} \arrow[l, shift left, "G"]
	\end{tikzcd}$,
	and suppose that $F$ is exact.
% segment-id: o014.li2.chapter3.k-injectives.l004
	\begin{compactitem}
% segment-id: o014.li2.chapter3.k-injectives.i007
		\item If $G$ is left adjoint to $F$, then
		$\cate{K}G:\cate{K}(\mathcal{B})\to\cate{K}(\mathcal{A})$
		carries K-projective complexes to K-projective complexes.
% segment-id: o014.li2.chapter3.k-injectives.i008
		\item If $G$ is right adjoint to $F$, then
		$\cate{K}G:\cate{K}(\mathcal{B})\to\cate{K}(\mathcal{A})$
		carries K-injective complexes to K-injective complexes.
	\end{compactitem}
\end{proposition}

% segment-id: o014.li2.chapter3.k-injectives.pf003
\begin{proof}
	By duality it suffices to consider the case in which $G$ is a left
	adjoint. In this case $G$ is automatically additive
	(Corollary~\ref{prop:automatic-additivity}). Let
	$P\in\Obj(\cate{C}(\mathcal{B}))$ be K-projective. For every acyclic
	complex $X\in\Obj(\cate{C}(\mathcal{A}))$,
	Proposition~\ref{prop:KF-adjoint} gives an isomorphism
	\[
	\Hom_{\cate{K}(\mathcal{A})}(\cate{K}G(P), X) \simeq
	\Hom_{\cate{K}(\mathcal{B})}(P, \cate{K}F(X)).
	\]
	Since $F$ is exact, $\cate{K}F(X)$ is acyclic, so the right-hand side is
	$0$. Apply Lemma~\ref{prop:K-injective-criterion}.
\end{proof}

% segment-id: o014.li2.chapter3.k-injectives.p005
The next task in this section is to investigate the existence of K-injective
or K-projective resolutions. The argument is somewhat intricate; we begin
with the K-injective case.

% segment-id: o014.li2.chapter3.k-injectives.le002
\begin{lemma}\label{prop:K-injective-projlimit-aux}
	Consider a sequence of complexes over $\mathcal{A}$,
	$\cdots\to I_2\to I_1\to I_0$. Suppose that:
% segment-id: o014.li2.chapter3.k-injectives.l005
	\begin{itemize}
% segment-id: o014.li2.chapter3.k-injectives.i009
		\item every $I_k$ is K-injective;
% segment-id: o014.li2.chapter3.k-injectives.i010
		\item for every $k$ and $n$, the morphism
		$I_{k+1}^n\to I_k^n$ is an epimorphism admitting a section
		(in other words, it has a right inverse; see the discussion after
		Proposition~\ref{prop:split-ses});
% segment-id: o014.li2.chapter3.k-injectives.i011
		\item for every $n$, the object
		$I^n:=\varprojlim_k I_k^n$ exists in $\mathcal{A}$.
	\end{itemize}
	Then the induced differentials
	$d_I^n:=\varprojlim_k d_{I_k}^n:I^n\to I^{n+1}$ make
	$I:=(I^n,d_I^n)_n$ a K-injective complex.
\end{lemma}

% segment-id: o014.li2.chapter3.k-injectives.pf004
\begin{proof}
	Let $X\in\Obj(\cate{C}(\mathcal{A}))$ be acyclic. For each
	$k\in\Z_{\geq0}$, take the following portion of the complex
	$\Hom^\bullet(X,I_k)$:
% segment-id: o014.li2.chapter3.k-injectives.e001
	\begin{equation}\label{eqn:K-injective-projlimit-aux}\begin{tikzcd}
		[row sep=small, column sep=tiny,
		every cell/.append style = {font = \small}]
		\displaystyle\prod_n \Hom\left( X^n, I_k^{n-2} \right) \arrow[r] &
		\displaystyle\prod_n \Hom\left( X^n, I_k^{n-1} \right) \arrow[r] &
		\displaystyle\prod_n \Hom\left( X^n, I_k^n \right) \arrow[r] &
		\displaystyle\prod_n \Hom\left( X^n, I_k^{n+1} \right) \\
		\Hom^{-2}\left( X, I_k \right) \arrow[equal, u] &
		\Hom^{-1}\left( X, I_k \right) \arrow[equal, u] &
		\Hom^0\left( X, I_k \right) \arrow[equal, u] &
		\Hom^1\left( X, I_k \right) \arrow[equal, u]
	\end{tikzcd}\end{equation}
	This sequence is exact because $I_k$ is K-injective. Now consider the
	homomorphism induced by $I_{k+1}\to I_k$,
% segment-id: o014.li2.chapter3.k-injectives.q001
	\[
	\Hom\left(X^n,I_{k+1}^{n-2}\right)\to
	\Hom\left(X^n,I_k^{n-2}\right).
	\]
	Since there is a section $I_k^{n-2}\to I_{k+1}^{n-2}$, this
	homomorphism is an epimorphism. The same remains true after taking
	$\prod_n$. Thus, as $k$ varies, the leftmost term in
	\eqref{eqn:K-injective-projlimit-aux} satisfies the Mittag--Leffler
	condition (Definition~\ref{def:ML}).

	Now let $k\in\Z_{\geq0}$ vary in
	\eqref{eqn:K-injective-projlimit-aux}.
	Corollary~\ref{prop:ML-4-terms} gives an exact sequence
% segment-id: o014.li2.chapter3.k-injectives.g004
	\[\begin{tikzcd}
		[column sep=small, every cell/.append style = {font = \small}]
		\displaystyle\varprojlim_k \displaystyle\prod_n
		\Hom\left( X^n, I_k^{n-1} \right) \arrow[r] &
		\displaystyle\varprojlim_k \displaystyle\prod_n
		\Hom\left( X^n, I_k^n \right) \arrow[r] &
		\displaystyle\varprojlim_k \displaystyle\prod_n
		\Hom\left( X^n, I_k^{n+1} \right)
	\end{tikzcd}\]
	Interchanging $\varprojlim_k$ and $\prod_n$ now gives an exact
	sequence\footnote{Translator's note (O014-C064): in the third term the
	source omits the subscript $k$ on $\varprojlim_k$. This sentence
	explicitly interchanges $\varprojlim_k$ with $\prod_n$, and the vertical
	isomorphism below identifies the limit with $I^{n+1}$, so the target
	restores that subscript.}
% segment-id: o014.li2.chapter3.k-injectives.g005
	\[\begin{tikzcd}
		[every cell/.append style = {font = \small}]
		\displaystyle\prod_n \varprojlim_k
		\Hom\left( X^n, I_k^{n-1} \right) \arrow[r] &
		\displaystyle\prod_n \varprojlim_k
		\Hom\left( X^n, I_k^n \right) \arrow[r] &
		\displaystyle\prod_n \varprojlim_k
		\Hom\left( X^n, I_k^{n+1} \right) \\
		\displaystyle\prod_n \Hom\left( X^n, I^{n-1} \right)
		\arrow[u, "\sim" sloped] &
		\displaystyle\prod_n \Hom\left( X^n, I^n \right)
		\arrow[u, "\sim" sloped] &
		\displaystyle\prod_n \Hom\left( X^n, I^{n+1} \right)
		\arrow[u, "\sim" sloped]
	\end{tikzcd}\]
	This too is a portion of $\Hom^\bullet(X,I)$, and its cohomology at the
	middle term is precisely
	$\Hom_{\cate{K}(\mathcal{A})}(X,I)$. Hence
	$\Hom_{\cate{K}(\mathcal{A})}(X,I)=0$. Apply
	Lemma~\ref{prop:K-injective-criterion} to finish the proof.
\end{proof}

% segment-id: o014.li2.chapter3.k-injectives.le003
\begin{lemma}\label{prop:K-injective-truncated}
	Suppose that $\mathcal{A}$ has enough injective objects. For every
	$A\in\Obj(\cate{C}(\mathcal{A}))$, there is a commutative diagram in
	$\cate{C}(\mathcal{A})$
% segment-id: o014.li2.chapter3.k-injectives.g006
	\[\begin{tikzcd}
		\cdots \arrow[r] & I_2 \arrow[r] & I_1 \arrow[r] & I_0 \\
		\cdots \arrow[r] & \tau^{\geq -2} A\arrow[u, "f_2"'] \arrow[r] &
		\tau^{\geq -1} A \arrow[r] \arrow[u, "f_1"'] &
		\tau^{\geq 0} A \arrow[u, "f_0"']
	\end{tikzcd}\]
	with $\tau^{\geq k}A$ and the arrows between them as in
	Definition~\ref{def:truncation-cplx}, such that:
% segment-id: o014.li2.chapter3.k-injectives.l006
	\begin{itemize}
% segment-id: o014.li2.chapter3.k-injectives.i012
		\item every $I_k$ is an object of
		$\cate{C}^+(\mathcal{A})$ consisting of injective objects;
% segment-id: o014.li2.chapter3.k-injectives.i013
		\item every vertical arrow $f_k$ is both a monomorphism and a
		quasi-isomorphism;
% segment-id: o014.li2.chapter3.k-injectives.i014
		\item for every $k$ and $n$, the morphism
		$I_{k+1}^n\to I_k^n$ is an epimorphism admitting a section.
	\end{itemize}
\end{lemma}

% segment-id: o014.li2.chapter3.k-injectives.pf005
\begin{proof}
	We construct the diagram from right to left. All injective resolutions
	discussed below are implicitly monomorphisms.

	First apply Theorem~\ref{prop:resolution-existence} to
	$\tau^{\geq0}A\in\Obj(\cate{C}^+(\mathcal{A}))$ to obtain an injective
	resolution $f_0:\tau^{\geq0}A\to I_0$.

	Next, the definitions directly show that
	$\tau^{\geq-k-1}A\to\tau^{\geq-k}A$ is an epimorphism in every degree,
	and hence an epimorphism in $\cate{C}^+(\mathcal{A})$. Suppose that the
	required quasi-isomorphisms
% segment-id: o014.li2.chapter3.k-injectives.q002
	\[
	\tau^{\geq0}A\xrightarrow{f_0}I_0,\quad\ldots,\quad
	\tau^{\geq-k}A\xrightarrow{f_k}I_k
	\]
	and the corresponding commutative diagram have been constructed. Choose
	an injective resolution $H\xrightarrow{g}J$ of
	$H:=\Ker[\tau^{\geq-k-1}A\to\tau^{\geq-k}A]$ by
	Theorem~\ref{prop:resolution-existence}. The
	Horseshoe Lemma, Proposition~\ref{prop:horseshoe}, then gives a
	commutative diagram with exact rows in $\cate{C}(\mathcal{A})$:
% segment-id: o014.li2.chapter3.k-injectives.g007
	\[\begin{tikzcd}
		0 \arrow[r] & J \arrow[r] & I_{k+1} \arrow[r] & I_k \arrow[r] & 0 \\
		0 \arrow[r] & H \arrow[r] \arrow[u, "g"] &
		\tau^{\geq-k-1}A \arrow[u, "{f_{k+1}}"] \arrow[r] &
		\tau^{\geq-k}A \arrow[r] \arrow[u, "f_k"'] & 0
	\end{tikzcd}\]
	such that $f_{k+1}$ is an injective resolution. Exactness of the first
	row is equivalent to exactness of
	$0\to J^n\to I_{k+1}^n\to I_k^n\to0$ for every $n$. Since all its terms
	are injective objects of $\mathcal{A}$,
	Lemma~\ref{prop:inj-proj-split} shows that
	$I_{k+1}^n\to I_k^n$ admits a section. This proves the result.
\end{proof}

% segment-id: o014.li2.chapter3.k-injectives.p006
The following discussion uses notation from \S\ref{sec:lim1}.

% segment-id: o014.li2.chapter3.k-injectives.p007
Take $A$ as above. As $k\geq0$ varies, the complexes
$\tau^{\geq-k}A$ form an object of
$\cate{InvSys}(\cate{C}(\mathcal{A}))$, which we denote by $\tau A$. The
definition of the truncation functors gives a family of canonical morphisms
$A\to\tau^{\geq-k}A$. Considering the complexes
$\tau^{\geq-k}A$ degree by degree shows that their $\varprojlim_k$ exists and
that these canonical morphisms induce an isomorphism
% segment-id: o014.li2.chapter3.k-injectives.e002
\begin{equation}\label{eqn:A-lim-truncation}
	A \rightiso \varprojlim_{k \geq 0} \left( \tau^{\geq -k} A \right).
\end{equation}

% segment-id: o014.li2.chapter3.k-injectives.p008
Now consider the $I_k$. Suppose that $\mathcal{A}$ has countable products;
then $\cate{C}(\mathcal{A})$ also has countable products, constructed degree
by degree. For $A\in\Obj(\cate{C}(\mathcal{A}))$, consider the diagram in
Lemma~\ref{prop:K-injective-truncated}. It follows that
% segment-id: o014.li2.chapter3.k-injectives.q003
\[
I:=\varprojlim_k I_k \quad \text{exists}.
\]

% segment-id: o014.li2.chapter3.k-injectives.p009
More concretely, Definition~\ref{def:Delta-diff} and the discussion following
it place $I$ and $A$ in exact sequences
% segment-id: o014.li2.chapter3.k-injectives.q004
\begin{gather*}
	0 \to I \to \prod_k I_k \xrightarrow{\Delta_I} \prod_k I_k, \\
	0 \to A \to \prod_k \tau^{\geq -k} A
	\xrightarrow{\Delta_{\tau A}} \prod_k \tau^{\geq -k} A,
\end{gather*}
where $\Delta_I$ and $\Delta_{\tau A}$ are as in that definition. We obtain
commutative diagrams
% segment-id: o014.li2.chapter3.k-injectives.e003
\begin{equation}\label{eqn:K-injective-diagonal}\begin{tikzcd}
	I \arrow[hookrightarrow, r] \arrow[d] & \displaystyle\prod_k I_k \\
	\Cone(\Delta_I)[-1] \arrow[ru, "{\beta(\Delta_I)[-1]}"'] &
\end{tikzcd} \quad \begin{tikzcd}
	A \arrow[hookrightarrow, r] \arrow[d] &
	\displaystyle\prod_k \tau^{\geq -k} A \\
	\Cone(\Delta_{\tau A})[-1]
	\arrow[ru, "{\beta(\Delta_{\tau A})[-1]}"']
\end{tikzcd}\end{equation}
The two vertical arrows are defined by inclusion into the first coordinate of
the mapping cone; the reader may check this directly. Another explanation,
though more elaborate than necessary, can be given using the universal
property of the homotopy kernel
(Proposition~\ref{prop:homotopy-kernel-cokernel}).

% segment-id: o014.li2.chapter3.k-injectives.p010
On the other hand, $(f_k)_{k\geq0}$ and the universal property of
$\varprojlim$ give a morphism
% segment-id: o014.li2.chapter3.k-injectives.q005
\[
f:A\simeq\varprojlim_k(\tau^{\geq-k}A)\to\varprojlim_k I_k=I,
\]
which makes the following diagram commute:
% segment-id: o014.li2.chapter3.k-injectives.e004
\begin{equation}\label{eqn:K-injective-A-I}\begin{tikzcd}
	I \arrow[r, "\text{as in \eqref{eqn:K-injective-diagonal}}"
	inner sep=0.4em] & \Cone\left(\Delta_I\right)[-1] \\
	A \arrow[r] \arrow[u, "f"] &
	\Cone\left( \Delta_{\tau A} \right)[-1]
	\arrow[u, "\text{induced by all the $f_k$}"'].
\end{tikzcd}\end{equation}
Lemmas~\ref{prop:K-injective-projlimit-aux} and
\ref{prop:K-injective-truncated} show that $I$ is a K-injective complex. If
$f$ is a quasi-isomorphism, then $f$ is a K-injective resolution of $A$.
This last step requires an additional hypothesis.

% segment-id: o014.li2.chapter3.k-injectives.le004
\begin{lemma}\label{prop:K-injective-equiv}
	In the situation above, suppose in addition that $\mathcal{A}$ has exact
	countable products (Convention~\ref{con:exact-product}). Then $f$ is a
	quasi-isomorphism if and only if
	$A\to\Cone(\Delta_{\tau A})[-1]$ is a quasi-isomorphism.
\end{lemma}

% segment-id: o014.li2.chapter3.k-injectives.pf006
\begin{proof}
	Consider the natural commutative diagram\footnote{Translator's note
	(O014-C065): in the source, the first lower horizontal arrow is labelled
	$\Delta_A$. The lower row belongs to the inverse system $\tau A$, and
	the right side of that row already uses $\alpha(\Delta_{\tau A})$ and
	$\Cone(\Delta_{\tau A})$; the target therefore restores the label
	$\Delta_{\tau A}$.}
% segment-id: o014.li2.chapter3.k-injectives.g008
	\[\begin{tikzcd}
		\prod_k I_k \arrow[r, "\Delta_I"] &
		\prod_k I_k \arrow[r, "{\alpha(\Delta_I)}"] &
		\Cone\left(\Delta_I\right) \\
		\prod_k \tau^{\geq-k}A \arrow[r, "\Delta_{\tau A}"']
		\arrow[u, "{\prod_k f_k}"] &
		\prod_k \tau^{\geq-k}A \arrow[u, "{\prod_k f_k}"]
		\arrow[r, "{\alpha(\Delta_{\tau A})}"'] &
		\Cone\left(\Delta_{\tau A}\right) \arrow[u]
	\end{tikzcd}\]
	Taking the long exact sequences of the mapping cones
	\eqref{eqn:cone-long-exact-sequence} gives a commutative diagram with
	exact rows:
% segment-id: o014.li2.chapter3.k-injectives.g009
	\[\begin{tikzcd}
		[column sep=tiny, every cell/.append style = {font = \small}]
		\cdots \Hm^n\left( \displaystyle\prod_k I_k \right) \arrow[r] &
		\Hm^n\left( \Cone(\Delta_I) \right) \arrow[r] &
		\Hm^{n+1}\left( \displaystyle\prod_k I_k \right)
		\arrow[r, "{\Hm^{n+1}(\Delta_I)}" inner sep=1em] &
		\Hm^{n+1}\left( \displaystyle\prod_k I_k \right) \\
		\cdots \Hm^n\left( \displaystyle\prod_k \tau^{\geq-k} A \right)
		\arrow[r] \arrow[u, "{\Hm^n(\prod_k f_k)}"] &
		\Hm^n\left( \Cone(\Delta_{\tau A}) \right)
		\arrow[r] \arrow[u] &
		\Hm^{n+1}\left(\displaystyle\prod_k \tau^{\geq-k} A \right)
		\arrow[u, "{\Hm^{n+1}(\prod_k f_k)}"]
		\arrow[r, "{\Hm^{n+1}(\Delta_{\tau A})}"' inner sep=1.5em] &
		\Hm^{n+1}\left( \displaystyle\prod_k \tau^{\geq-k} A \right)
		\arrow[u, "{\Hm^{n+1}(\prod_k f_k)}"'] .
	\end{tikzcd}\]
	Every $f_k$ is a quasi-isomorphism, and exactness of countable products
	ensures that $\prod_k f_k$ remains a quasi-isomorphism. Applying
	Proposition~\ref{prop:5-lemma} shows that
	$\Cone(\Delta_{\tau A})\to\Cone(\Delta_I)$\footnote{Translator's note
	(O014-C066): the source writes $\Delta_{\tau_A}$ at this occurrence,
	although the inverse-system object that was defined is $\tau A$ and all
	surrounding notation uses $\Delta_{\tau A}$. The target makes this
	consistent with the defined notation.} is also a quasi-isomorphism.

	Next we prove that $I\to\Cone(\Delta_I)[-1]$ is a quasi-isomorphism.
	First, $\Delta_I$ is an epimorphism; this can be checked degree by degree.
	Indeed,
	$\Delta_{I^n}:\prod_k I_k^n\to\prod_k I_k^n$ is an epimorphism because
	$I_{k+1}^n\to I_k^n$ admits a section; see
	Example~\ref{eg:lim1-section}. Compare the short exact sequence
	$0\to I\to\prod_k I_k\xrightarrow{\Delta_I}\prod_k I_k\to0$ with
	Lemma~\ref{prop:cone-qis}. It follows that the previously defined
	morphism $I\to\Cone(\Delta_I)[-1]$ is indeed a quasi-isomorphism.

	Insert these results into the commutative diagram
	\eqref{eqn:K-injective-A-I} and then take cohomology. This gives the
	desired necessary and sufficient condition.
\end{proof}

% segment-id: o014.li2.chapter3.k-injectives.th002
\begin{theorem}\label{prop:K-injective-resolution}
	An abelian category $\mathcal{A}$ satisfying the following two
	conditions has enough K-injective complexes:
% segment-id: o014.li2.chapter3.k-injectives.l007
	\begin{compactitem}
% segment-id: o014.li2.chapter3.k-injectives.i015
		\item $\mathcal{A}$ has exact countable products;
% segment-id: o014.li2.chapter3.k-injectives.i016
		\item $\mathcal{A}$ has enough injective objects.
	\end{compactitem}
\end{theorem}

% segment-id: o014.li2.chapter3.k-injectives.pf007
\begin{proof}
	Lemma~\ref{prop:K-injective-equiv} shows that it suffices to prove that
	$A\to\Cone(\Delta_{\tau A})[-1]$ is a quasi-isomorphism for every
	$A\in\Obj(\cate{C}(\mathcal{A}))$. Let $n\in\Z$. Since countable
	products are exact, Lemma~\ref{prop:truncation-ses} gives canonical
	isomorphisms
% segment-id: o014.li2.chapter3.k-injectives.q006
	\[
	\Hm^n\left( \prod_k \tau^{\geq -k} A \right)
	\simeq \prod_k \Hm^n\left( \tau^{\geq -k} A \right)
	\simeq \prod_{k \geq -n} \Hm^n(A).
	\]
	If elementwise reasoning is available in $\mathcal{A}$, then after
	taking $\Hm^n$, the morphism $\Delta_{\tau A}$ has the form
% segment-id: o014.li2.chapter3.k-injectives.g010
	\[\begin{tikzcd}[row sep=tiny]
		\prod_{k \geq -n} \Hm^n(A) \arrow[r, "\nabla^n"] &
		\prod_{k \geq -n} \Hm^n(A) \\
		(a_k)_{k \geq -n} \arrow[mapsto, r] &
		(a_{k+1} - a_k)_{k \geq -n}.
	\end{tikzcd}\]
	Clearly $\nabla^n$ is an epimorphism (a section can be written down
	explicitly), while
	$\Ker(\nabla^n)=\{(a)_{k\geq-n}:a\in\Hm^n(A)\}$ is the diagonal
	subobject. This argument readily translates to a general category
	$\mathcal{A}$. Now consider the long exact sequence of the mapping cone
% segment-id: o014.li2.chapter3.k-injectives.q007
	\[
	\cdots \to \Hm^n \Cone(\Delta_{\tau A})[-1]
	\to \prod_{k \geq -n} \Hm^n(A)
	\xrightarrow[\text{epi}]{\nabla^n}
	\prod_{k \geq -n} \Hm^n(A)
	\to \Hm^n \Cone(\Delta_{\tau A}) \cdots
	\]
	This sequence identifies the morphism
	$\Hm^n\Cone(\Delta_{\tau A})[-1]\to
	\prod_{k\geq-n}\Hm^n(A)$ with the inclusion of
	$\Ker(\nabla^n)$. Moreover, \eqref{eqn:K-injective-diagonal} gives a
	commutative diagram
% segment-id: o014.li2.chapter3.k-injectives.g011
	\[\begin{tikzcd}
		\Hm^n(A) \arrow[d]
		\arrow[hookrightarrow, r, "\text{diagonal morphism}"] &
		\prod_{k \geq -n} \Hm^n(A) \\
		\Hm^n {\Cone(\Delta_{\tau A})[-1]} \arrow[hookrightarrow, ru] &
	\end{tikzcd}\]
	Thus, for every $n$, the morphism
	$\Hm^n(A)\to\Hm^n(\Cone(\Delta_{\tau A})[-1])$ is an isomorphism.
\end{proof}

% segment-id: o014.li2.chapter3.k-injectives.p011
Reversing the arrows by means of Proposition~\ref{prop:K-sigma-duality}
gives the dual version of Theorem~\ref{prop:K-injective-resolution}.

% segment-id: o014.li2.chapter3.k-injectives.th003
\begin{theorem}\label{prop:K-projective-resolution}
	An abelian category $\mathcal{A}$ satisfying the following two
	conditions has enough K-projective complexes:
% segment-id: o014.li2.chapter3.k-injectives.l008
	\begin{compactitem}
% segment-id: o014.li2.chapter3.k-injectives.i017
		\item $\mathcal{A}$ has exact countable coproducts;
% segment-id: o014.li2.chapter3.k-injectives.i018
		\item $\mathcal{A}$ has enough projective objects.
	\end{compactitem}
\end{theorem}

% segment-id: o014.li2.chapter3.k-injectives.co001
\begin{corollary}\label{prop:Grothendieck-K-projective}
	If $\mathcal{A}$ is a Grothendieck category with enough projective
	objects (Definition~\ref{def:Grothendieck-cat}), then $\mathcal{A}$ also
	has enough K-projective complexes.
\end{corollary}

% segment-id: o014.li2.chapter3.k-injectives.pf008
\begin{proof}
	By definition, a Grothendieck category $\mathcal{A}$ has countable
	coproducts. Proposition~\ref{prop:Grothendieck-cat-coprod-exact} says
	that those countable coproducts are exact. Apply
	Theorem~\ref{prop:K-projective-resolution}.
\end{proof}

% segment-id: o014.li2.chapter3.k-injectives.ex003
\begin{example}\label{eg:Mod-K-injectives}
	Let $R$ be a ring. Theorem~\ref{prop:K-injective-resolution} and
	Corollary~\ref{prop:Grothendieck-K-projective} show that
	$R\dcate{Mod}$ has enough K-injective and K-projective complexes.
\end{example}

% segment-id: o014.li2.chapter3.k-injectives.p012
In fact, every complex in a Grothendieck category has a K-injective
resolution, and these K-injective resolutions can be made functorial, as in
Theorem~\ref{prop:Grothendieck-injective}. This is the main theorem of
\cite{Serp03}; see also \cite[Tag 079P]{stacks}. That result covers many
situations arising in geometry that are not covered by
Theorem~\ref{prop:K-injective-resolution}.
